\newcommand{\cano}{K} % canonical divisor
\newcommand{\DValo}{\DVal^\circ}
\newcommand{\lct}{\mathrm{lct}}

In this section, given a projective variety $X$ with a canonical divisor $\cano_X$ (which always exists as $X$ is assumed to be normal). The variety $X$ is called \emph{$\mathbb{Q}$-Gorenstein} if $\cano_X$ is $\mathbb{Q}$-Cartier.

Assume $X$ is $\mathbb{Q}$-Gorenstein. For every prime divisor $E$ on a birational model $\pi: Y \to X$, we can define the \emph{log discrepancy} of divisorial valuations $c \: \ord_E$ on $X$ as follows
\begin{equation}
	A_X(c \: \ord_E) := c \cdot (1 + \ord_E(\cano_Y - \pi^* \cano_X)) ,
\end{equation}
where $\cano_Y$ is a canonical divisor of $Y$.

\begin{remark}
	For the definition of log discrepancy for general $\mathbb{R}$-valuations, see \cite[\S5]{JM_val_asym_inva_seq_i}.
\end{remark}

A $\mathbb{Q}$-Gorenstein variety $X$ is called \emph{Kawamata log terminal} (in short, \emph{klt}) if $A_X(\ord_E) > 0$ for every divisor $E$ over $X$.

\begin{definition}
	Given a klt variety $X$, for a non-zero effective $\mathbb{Q}$-Cartier $\mathbb{Q}$-divisor $D$ on $X$, the \emph{log canonical threshold} of $D$ is defined by,
	\begin{equation*}
		\lct(D) := \lct(X; D) = \inf_{v \in \DValo_{X}} {A_X(v) \over v(D)} ;
	\end{equation*}
	where $\DValo_{X} : = \{v \in \DVal_X \mid v \neq 0 \}$.
\end{definition}
